\chapter{When Markets Break Together}\label{ch:m3:when-markets-break-together}

In the middle of March 2020, as the pandemic closed the world's economies, equities fell, and so did the
assets that are supposed to rise when equities fall. Between 9 and 18 March the yield on ten-year US
Treasuries, the world's safest asset, rose from 0.54\% to 1.18\% as holders sold them; the broad dollar
index rose more than 7\%; on 16 March the S\&P 500 lost 12\%, its worst day since 1987. Investors were
not choosing between risky and safe assets any more; they were selling whatever they could to raise cash.
Earlier chapters described each market's own crises: power prices in 2022, nickel on the London Metal
Exchange, crypto \omterm{def:m3:margin-liquidation-and-loss-allocation:cascade}{liquidation cascades}, FTX. This chapter is about the days when the crises are the same
crisis, and the connections between markets that are invisible in calm, common holders, common
financing, common margin rules, become the only thing that matters. It sets out the mechanics, goes
through 2008, March 2020 and 2022 from the point of view of the desks in this book, and builds a
simulator of the spiral.

\section{The mechanics of joint stress: fire sales, loss and margin spirals}

\begin{definition}[Fire sale]\label{def:m3:when-markets-break-together:fire}
A \emph{fire sale}\index{fire sale} is a sale forced on a holder by a need for cash or a margin call, at a price below what the asset is
worth to buyers who are not themselves constrained, because the natural buyers of the asset are selling too.
\end{definition}

\begin{definition}[Loss spiral, margin spiral]\label{def:m3:when-markets-break-together:spirals}
A \emph{loss spiral}\index{loss spiral} is the feedback in which a leveraged holder's losses force it to sell, its sales lower prices, and the
lower prices inflict further losses on it and on other holders of the same assets. A \emph{margin spiral}\index{margin spiral} is the feedback
in which rising volatility raises margin requirements, the higher margins force deleveraging, and the forced sales raise volatility further.
\end{definition}

Brunnermeier and Pedersen gave the two spirals their names in a model of how funding liquidity and market liquidity feed each other, the liquidity
spiral of One Quant Book 1, chapter 31. Two properties turn them into a cross-market event. Leveraged holders own many assets, so a loss in one forces
sales of all of them; and margins and haircuts are set on each asset's own recent volatility, so they rise everywhere at once. The asset that is sold
is not the one that fell, but the one that can be sold.

\begin{omfigure}
\centering
\begin{tikzpicture}[font=\small,
  s/.style={draw,thick,rounded corners=2pt,align=center,minimum height=0.9cm,text width=2.6cm},
  arr/.style={-{Stealth[length=2.2mm]},thick}]
\node[s,draw=omThm,fill=omThm!8] (sh) at (0,0) {shock: one\\asset falls};
\node[s,draw=omDef,fill=omDef!8] (eq) at (3.8,1.1) {holders' equity\\falls};
\node[s,draw=omExo,fill=omExo!8] (mg) at (3.8,-1.1) {volatility and\\margins rise};
\node[s,draw=omProp,fill=omProp!8] (fs) at (7.6,0) {forced sales of\\\emph{all} assets held};
\node[s,draw=gray,fill=gray!8] (px) at (11.2,0) {prices fall\\everywhere};
\draw[arr] (sh) -- (eq); \draw[arr] (sh) -- (mg); \draw[arr] (eq) -- (fs); \draw[arr] (mg) -- (fs); \draw[arr] (fs) -- (px);
\draw[arr,dashed] (px.north) to[bend right=30] node[above,font=\footnotesize] {loss spiral} (eq.north);
\draw[arr,dashed] (px.south) to[bend left=30] node[below,font=\footnotesize] {margin spiral} (mg.south);
\end{tikzpicture}

\omcaption{The two spirals. A shock to one asset cuts leveraged holders' equity and raises measured volatility and margins; both force sales of
everything the holders own; the sales lower prices and raise volatility everywhere, feeding back into both. Schematic.}\label{fig:m3:when-markets-break-together:spirals}
\end{omfigure}

\begin{proposition}[The forced-sale multiplier without impact]\label{prop:m3:when-markets-break-together:multiplier}
A holder whose equity falls short of its margin requirement by $S$, on positions with margin rate $m$, must sell $S/m$ of them to restore its margin,
if prices do not move: each dollar sold releases $m$ of requirement and leaves equity unchanged. With price impact, each sale lowers the prices of the
positions kept and the equity with them, and the forced sales per dollar of initial shortfall exceed $1/m$.
\end{proposition}

\begin{proof}
Selling a dollar of position at the market price changes equity by nothing and the requirement by $m$; restoring $S$ takes $S/m$. With impact the sold
positions move prices against the remaining ones, creating a new shortfall that must be met by further sales.
\end{proof}

\begin{definition}[Flight to quality, correlation breakdown]\label{def:m3:when-markets-break-together:flight}
A \emph{flight to quality}\index{flight to quality} is a shift of investors out of risky assets into those perceived as safest, typically raising the
prices of government bonds while other assets fall. A \emph{correlation breakdown}\index{correlation breakdown} is a sudden change in the correlations
between assets in a stress, usually towards one as forced selling moves them together, which invalidates hedges and risk estimates calibrated in calm.
\end{definition}

\section{2008}

The crisis of 2008 is the template: leveraged institutions financed short term against securities whose values and haircuts moved together, and the
failure of one of them, Lehman Brothers, on 15 September 2008, turned doubts about each into a run on all. The clearing houses of One Quant Book 1 held:
LCH's swap service managed Lehman's portfolio of 66\,390 trades with a notional of USD~9 trillion within the margin it held, by hedging and auctioning it
between 24 September and 3 October, without using its default fund. The losses came elsewhere, in the uncleared and the financed: short-term funding that stopped rolling
and haircuts that rose on collateral of every kind. The
lessons that shaped the post-2008 system, central clearing, margin for uncleared derivatives, bank liquidity rules, are the reason the next crisis looked
different.

\section{March 2020}

\begin{dated}{2026-09}{The March 2020 turmoil}\label{dat:m3:when-markets-break-together:2020}
The Financial Stability Board's review distinguishes a flight to safety from late February to early March 2020, in which investors sold risky assets and
bought safe ones, and a ``dash for cash'' in mid-March, in which they sold risky and relatively safe assets alike to obtain cash; equities and government
bonds then fell together, volatility rose even for government bonds and gold, and correlations broke down. On 16 March the S\&P 500 lost 12\%, its worst
day since 1987, and the VIX reached its all-time peak of 83. Market dysfunction in US Treasuries was exacerbated by sales by some leveraged non-bank
investors and foreign holders. The Basel Committee, CPMI and IOSCO reported that initial margin requirements at central counterparties rose by roughly USD~300
billion over March 2020, with a peak variation margin call of USD~140 billion on 9 March.
\end{dated}

\begin{omfigure}
\begin{tikzpicture}
\begin{axis}[width=11cm,height=5cm,axis y line*=left,
  xlabel={trading day from 3 February 2020},ylabel={10-year Treasury yield, \%},
  tick label style={font=\small},label style={font=\small},
  xmin=0,xmax=61,ymin=0.4,ymax=1.7,grid=major,grid style={gray!25},
  table/col sep=comma]
\fill[omExo!15] (axis cs:24,0.4) rectangle (axis cs:34,1.7);
\addplot[omThm,very thick] table[x=day,y=ust10y]{figdata/markets-3/28-when-markets-break-together/march2020.csv};\label{plot:m3:ust10}
\node[font=\footnotesize,anchor=north,text=omExo] at (axis cs:29,1.68) {dash for cash};
\end{axis}
\begin{axis}[width=11cm,height=5cm,axis y line*=right,axis x line=none,
  ylabel={broad dollar index},ylabel near ticks,
  tick label style={font=\small},label style={font=\small},
  xmin=0,xmax=61,ymin=112,ymax=128,
  legend columns=2,legend style={font=\footnotesize,at={(0.5,-0.3)},anchor=north,draw=none,fill=none,/tikz/every even column/.append style={column sep=8pt}},
  table/col sep=comma]
\addlegendimage{/pgfplots/refstyle=plot:m3:ust10}\addlegendentry{10-year yield (left)}
\addplot[omDef,thick,dashed] table[x=day,y=usd]{figdata/markets-3/28-when-markets-break-together/march2020.csv};\addlegendentry{dollar index (right)}
\end{axis}
\end{tikzpicture}

\omcaption{Safe assets in March 2020: the ten-year Treasury yield fell to 0.54\% on 9 March in the flight to safety, then rose to 1.18\% by 18 March as
Treasuries were sold for cash (shaded, 9 to 23 March), while the broad dollar index rose 7.6\% from 9 to 23 March. Data: FRED series DGS10 and DTWEXBGS (Board
of Governors of the Federal Reserve System).}\label{fig:m3:when-markets-break-together:march2020}
\end{omfigure}

The desks of this book saw the same week through different windows. The rates desk saw Treasuries sold with equities and the futures basis blow out
(One Quant Book 2). The commodity desk saw oil collapse for reasons of its own (the storage chapters of this book, and the negative WTI settlement of the
next month) while margins on everything rose. A crypto desk saw bitcoin sold with everything else, whatever its reputation as a hedge: its Coinbase price,
as FRED records it, fell from 7\,957 dollars on 11 March to 4\,980 on 12 March, 37\% in a day. Every desk saw its margins rise the same week, which is the \omterm{def:m3:when-markets-break-together:spirals}{margin spiral} seen from the inside.

\section{2022: energy margins, nickel, gilts and crypto}

2022 had no single crash but a sequence of spirals in different markets, each driven by the same mechanics in a different place. European gas and power prices
rose so far (\cref{ch:m3:natural-gas-and-lng}) that the cash margin on futures hedges became, for firms hedging physical positions, the liquidity risk
described in \cref{ch:m3:getting-access-commodities-and-power}. The ECB found initial margin on commodity portfolios about twice as high by mid-2022 as in
late 2021, commodities making up about 2\% of the gross notional of euro-area derivatives but over 20\% of the initial margin posted, and four banks
clearing about 85\% of exchange-traded energy positions; the gas importer Uniper had drawn the whole of a EUR~9 billion credit line from the German state
bank KfW by 29 August, and asked for 4 billion more, as gas curtailments and prices raised its margining needs.
In March, nickel on the London Metal Exchange rose from about 30\,000 to over 100\,000 dollars a tonne in a day and a morning, the day's margin calls reached
about USD~7 billion and three clearing members failed to pay on time, and the exchange suspended trading and cancelled a morning's trades
(\cref{ch:m3:metals}). In September, UK pension funds' liability-driven investment strategies met margin calls on their gilt positions as long yields rose
faster than ever, sold gilts to pay them, and pushed yields higher, until the Bank of England bought long gilts (One Quant Book 2, chapter 7). In crypto, the
\omterm{def:m3:defi-credit-and-leverage:algo}{algorithmic stablecoin} UST lost its peg in May and fell close to zero (\cref{ch:m3:defi-credit-and-leverage}); Three Arrows Capital, a crypto fund reported to manage over
USD~3 billion that had borrowed coins and dollars from several lenders, received notices of default from them as prices fell and was put into liquidation in
the British Virgin Islands on 27 June; and FTX failed in November
(\cref{ch:m3:centralised-exchanges}).

\section{What each desk saw}

The same mechanics look different from each desk. For the trading firm the lessons are practical. Collateral is the constraint: in every episode the
losers were not wrong about value but short of cash when margins rose, so the firm's liquidity plan must assume that margins rise on all its venues at
once. Hedges fail when correlations break: a hedge estimated in calm (bonds against equities, bitcoin against equities, one venue against another) may move
the same way in a dash for cash. Common holders make distant markets neighbours: the firm must know who else holds what it holds, with what leverage.
And stress tests (One Quant Book 6) must include the joint scenario, not only each market's worst day. The simulator of this chapter makes the point with
numbers.

\begin{omfigure}
\begin{tikzpicture}
\begin{axis}[width=11cm,height=5cm,
  xlabel={day},ylabel={price},
  tick label style={font=\small},label style={font=\small},
  xmin=0,xmax=39,ymin=70,ymax=110,grid=major,grid style={gray!25},
  legend columns=3,legend style={font=\footnotesize,at={(0.5,-0.3)},anchor=north,draw=none,fill=none,/tikz/every even column/.append style={column sep=8pt}},
  table/col sep=comma]
\addplot[omThm,very thick] table[x=day,y=a1]{figdata/markets-3/28-when-markets-break-together/paths.csv};\addlegendentry{asset 1 (shocked)}
\addplot[omDef,very thick,dashed] table[x=day,y=a2]{figdata/markets-3/28-when-markets-break-together/paths.csv};\addlegendentry{asset 2}
\addplot[omExo,very thick,dotted] table[x=day,y=a3]{figdata/markets-3/28-when-markets-break-together/paths.csv};\addlegendentry{asset 3}
\draw[gray,thick] (axis cs:10,70) -- (axis cs:10,110);
\end{axis}
\end{tikzpicture}

\omcaption{The tutorial's \omterm{def:m3:when-markets-break-together:spirals}{margin spiral}: 40 holders with 8 times leverage across three assets; on day 10 asset 1 falls 10\%. Rising margins and falling
equity force sales of all three, and assets 2 and 3, untouched by the shock, fall with it; the correlation of assets 1 and 2 rises from 0.30 in the ten days
before to 0.99 in the ten days from the shock, which the joint fall on the shock day dominates. Synthetic. Data: the chapter's tutorial.}\label{fig:m3:when-markets-break-together:paths}
\end{omfigure}

\begin{table}[ht]
\centering
\small
\begin{tabular}{@{}lrrr@{}}
\toprule
Case & forced sales, USD bn & multiplier & correlation after \\
\midrule
Base case & 223.6 & 8.88 & 0.99 \\
Eight steps a day & 225.6 & 8.73 & 0.99 \\
No \omterm{def:m3:when-markets-break-together:spirals}{margin spiral} (fixed margins) & 0.0 & 0.00 & 0.11 \\
No \omterm{def:m3:when-markets-break-together:spirals}{loss spiral} (equity not marked) & 121.9 & 4.94 & 0.98 \\
No price impact & 124.4 & 4.94 & 0.11 \\
One asset per holder & 183.8 & 4.67 & 0.04 \\
\bottomrule
\end{tabular}
\caption{The spiral and its ablations: forced sales over the run, forced sales per dollar of the margin shortfall the shock created (the forced-sale
multiplier), and the correlation of assets 1 and 2 over the ten days after the shock. Without price impact the multiplier is the $1/m$ of
\cref{prop:m3:when-markets-break-together:multiplier}; with it, it nearly doubles; a finer time step changes it little; with fixed margins, this shock is
absorbed without forced sales.}\label{tab:m3:when-markets-break-together:ablations}
\end{table}

\section{Tutorial: a margin spiral across three assets}\label{tut:m3:when-markets-break-together:spiral}

\begin{tutorial}
\textbf{Goal.} Simulate a \omterm{def:m3:when-markets-break-together:spirals}{margin spiral} across three assets held by leveraged investors, with margins rising with volatility and forced sales moving prices,
measure how correlations rise in the stress, check the result under a finer time step, and ablate each mechanism. \textbf{End state:}
\cref{fig:m3:when-markets-break-together:paths}, \cref{tab:m3:when-markets-break-together:ablations} and the numbers of the weekend problem.
\begin{tutsteps}
\item \textbf{Deleveraging.} Each holder short of margin sells the same fraction of everything; sales move prices.
\omcode{code/firm/marginspiral/firm_marginspiral.py}{100}{119}{Forced sales and their price impact.}\label{lst:m3:when-markets-break-together:delever}
\item \textbf{Stability and ablations.} \texttt{ablation\_table()} runs the base case, eight steps a day, and each mechanism switched off.
\item \textbf{Run} \texttt{correlations()} and \texttt{fig\_break.py}.
\end{tutsteps}
\textbf{What to change next.} Give holders different leverage and see who defaults first; add a cash buffer and find the size at which the spiral stops;
let margins respond with a lag, as a clearing house's model with a look-back would.
\end{tutorial}

\section{Build: the margin-spiral simulator}\label{bld:m3:when-markets-break-together:marginspiral}

\begin{build}
\bfield{Purpose} The miniature firm's stress tests must include the joint scenario in which its own margins rise on every venue while the assets it holds fall
together; the simulator produces those scenarios and measures how much forced selling a shock creates.
\bfield{Interface} \texttt{Params} (volatilities, correlation, base margins, impact, holders, leverage, days, shock, EWMA decay, substeps, switches for the \omterm{def:m3:when-markets-break-together:spirals}{margin
spiral}, the \omterm{def:m3:when-markets-break-together:spirals}{loss spiral}, cross-holding and forced selling); \texttt{simulate(params)} returning prices, returns, forced sales, the initial shortfall and
defaults; \texttt{forced\_sale\_multiplier}; \texttt{correlation(result, a, b, days)}.
\bfield{Rules} Deleveraging iterated to a fixed point within each step, so that results do not depend on the step size; log price impact, so prices stay
positive; every mechanism can be switched off.
\bfield{Acceptance tests} \texttt{code/firm/marginspiral/tests/}: the multiplier equals $1/m$ without impact and \omterm{def:m3:when-markets-break-together:spirals}{loss spiral}; stability under a finer time step;
each ablation lowers forced sales or correlation; no forced sales when forced selling is off.
\bfield{Stretch} Heterogeneous holders and a cash buffer; lagged margin models; a second shock; calibration of impact to a market's depth.
\end{build}

\begin{omsources}
\item M.~K.~Brunnermeier and L.~H.~Pedersen, ``Market Liquidity and Funding Liquidity'', \emph{Review of Financial Studies} 22(6), 2009.
\item Financial Stability Board, \emph{Holistic Review of the March Market Turmoil}, 17 November 2020.
\item BCBS--CPMI--IOSCO, press release on margining practices, 29 September 2022; LCH.Clearnet press release on the Lehman default, 8 October 2008 (see One
Quant Book 1's ledgers).
\item FRED, series DGS10, DTWEXBGS and CBBTCUSD; chapters 8, 13, 15 and 23 of this book and One Quant Book 2, chapter 7, for 2022.
\item ECB, \emph{Financial Stability Review}, November 2022, special feature ``Financial stability risks from energy derivatives markets''; Fortum, press
release on Uniper's KfW facility, 29 August 2022.
\item In re Three Arrows Capital Ltd., No.\ 22-10920 (Bankr.\ S.D.N.Y.), verified petition and declaration of the joint liquidators, 1 July 2022.
\end{omsources}

\section{Exercises}

\begin{exercise}[$\star$]\label{exo:m3:when-markets-break-together:1}
A holder is short of margin by USD~10 million on positions margined at 12.5\%. How much must it sell if prices do not move?
\end{exercise}

\begin{exercise}[$\star$]\label{exo:m3:when-markets-break-together:2}
Distinguish a \omterm{def:m3:when-markets-break-together:flight}{flight to quality} from a dash for cash, and say what each does to government bond prices.
\end{exercise}

\begin{exercise}[$\star$]\label{exo:m3:when-markets-break-together:3}
Why can an asset regarded as a hedge fall with equities in a dash for cash?
\end{exercise}

\begin{exercise}[$\star\star$]\label{exo:m3:when-markets-break-together:4}
From \cref{tab:m3:when-markets-break-together:ablations}, how much of the base case's forced selling is due to price impact and the \omterm{def:m3:when-markets-break-together:spirals}{loss spiral} together?
\end{exercise}

\begin{exercise}[$\star\star$]\label{exo:m3:when-markets-break-together:5}
The ten-year yield rose from 0.54\% to 1.18\% between 9 and 18 March 2020. With a modified duration of 9, what did a ten-year Treasury lose?
\end{exercise}

\begin{exercise}[$\star\star$]\label{exo:m3:when-markets-break-together:6}
Why does a \omterm{def:m3:when-markets-break-together:flight}{correlation breakdown} invalidate a value-at-risk estimated on calm data?
\end{exercise}

\begin{exercise}[$\star\star\star$]\label{exo:m3:when-markets-break-together:7}
\emph{Coding.} With \texttt{run}, halve the leverage. What happens to forced sales and the multiplier?
\end{exercise}

\begin{exercise}[$\star\star\star$]\label{exo:m3:when-markets-break-together:8}
\emph{Find the flaw.} ``Our portfolio is diversified across rates, commodities and crypto, so a crisis in one will not hit the others.''
\end{exercise}

\section{Problem: March 2020 Across the Desks}

\begin{problem}[{Weekend problem --- how much selling does a shock force?}]\label{pb:m3:when-markets-break-together:1}
Use the tutorial's simulator: 40 holders, each with USD~1 billion of equity and 8 times leverage spread over three assets, margins rising with EWMA volatility, and
a 10\% shock to asset 1 on day 10.

\medskip\noindent\textbf{Part I --- The shock.}
\begin{enumerate}
\item What margin shortfall does the shock create across the holders?
\item How much do they sell by force over the run?
\item What is the forced-sale multiplier?
\item What would the multiplier be without price impact, and why?
\item What happens to assets 2 and 3, which the shock did not touch?
\end{enumerate}

\medskip\noindent\textbf{Part II --- The mechanisms.}
\begin{enumerate}[resume]
\item What does the result become with eight steps a day, and why does that matter?
\item What happens with fixed margins?
\item And if each holder owns only one asset?
\item How does the correlation between assets 1 and 2 change after the shock?
\item Which mechanism does most of the amplification?
\end{enumerate}

\medskip\noindent\textbf{Part III --- The history.}
\begin{enumerate}[resume]
\item What did the rates desk see in March 2020?
\item What did clearing houses do to margins that month?
\item Why did LCH's handling of Lehman in 2008 not prevent the crisis?
\item Which 2022 episode most resembles the simulator, and why?
\item What did central banks do in each episode that the simulator leaves out?
\end{enumerate}

\medskip\noindent\textbf{Part IV --- Judgement.}
\begin{enumerate}[resume]
\item Should clearing houses raise margins in a crisis?
\item How should a trading firm size its liquidity buffer for joint stress?
\item What would you monitor to see a spiral coming?
\item State the \emph{named result}: the forced-sale multiplier of the spiral for the stated leverage and impact, and without impact.
\item In one sentence: why do markets break together?
\end{enumerate}
\end{problem}

\section{Interview questions}

\begin{interviewq}[$\star$ \iqroles{risk}]\label{iq:m3:when-markets-break-together:1}
What are \omterm{def:m3:when-markets-break-together:spirals}{loss spirals} and \omterm{def:m3:when-markets-break-together:spirals}{margin spirals}?
\end{interviewq}

\begin{interviewq}[$\star$ \iqroles{trader}]\label{iq:m3:when-markets-break-together:2}
Why did Treasuries fall in the middle of March 2020?
\end{interviewq}

\begin{interviewq}[$\star\star$ \iqroles{researcher}]\label{iq:m3:when-markets-break-together:3}
How would you estimate correlations for a stress scenario?
\end{interviewq}

\begin{interviewq}[$\star\star$ \iqroles{risk}]\label{iq:m3:when-markets-break-together:4}
Design a liquidity stress test for a firm trading futures, crypto perpetuals and FX forwards.
\end{interviewq}

\begin{interviewq}[$\star\star$ \iqroles{trader, risk}]\label{iq:m3:when-markets-break-together:5}
Your hedge fund client is long Treasuries and short futures with 50 times leverage. What happens to it in a dash for cash?
\end{interviewq}

\begin{interviewq}[$\star\star\star$ \iqroles{researcher, developer}]\label{iq:m3:when-markets-break-together:6}
Build a simulator of a \omterm{def:m3:when-markets-break-together:spirals}{margin spiral} and explain how you would make sure its results are not an artefact of its time step.
\end{interviewq}
